\chapter{Análise dos Resultados}
\label{chap:resultados}

\section{Aquisição e Carga de Dados}

Os 16 sujeitos do dataset IEEE Multimodal Emotion Recognition foram processados integralmente pelo estágio inicial do pipeline de aquisição. A Tabela \ref{tab:resumo_aquisicao} resume os parâmetros técnicos e temporais obtidos para as modalidades de biossinais coletadas.

\begin{table}[htbp]
    \centering
    \caption{Sumário técnico dos biossinais adquiridos}
    \label{tab:resumo_aquisicao}
    \begin{tabular}{lccc}
        \toprule
        \textbf{Modalidade} & \textbf{Canais} & \textbf{Freq. Amostragem (Hz)} & \textbf{Duração (s)} \\
        \midrule
        EEG & 8 & 512 & 60 -- 145 \\
        ECG & 1 & 250 & 60 -- 260 \\
        EMG & 1 & 250 & 60 -- 260 \\
        fNIRS & 2 & 16 & 60 -- 145 \\
        \bottomrule
    \end{tabular}
\end{table}

\section{Validação de Nyquist}

Conforme planejado, todas as modalidades apresentam conformidade estrita com o teorema de amostragem de Nyquist, demonstrando razões entre a taxa de amostragem e a frequência máxima superiores ao limite teórico de 2,00. Esta conformidade assegura que os sinais foram amostrados adequadamente e podem ser submetidos a filtragens digitais e transformadas espectrais sem distorções de \textit{aliasing} na banda de interesse, conforme detalhado na Tabela \ref{tab:nyquist}.

\begin{table}[htbp]
    \centering
    \caption{Parâmetros de validação conforme o teorema de Nyquist}
    \label{tab:nyquist}
    \begin{tabular}{lcccc}
        \toprule
        \textbf{Modalidade} & $f_s$ (Hz) & $f_{max}$ (Hz) & $f_s / 2f_{max}$ & \textbf{Conforme} \\
        \midrule
        EEG & 512 & 50 & 5,12 & \checkmark \\
        ECG & 250 & 40 & 3,13 & \checkmark \\
        EMG & 250 & 100 & 1,25 & \checkmark \\
        fNIRS & 16 & 2 & 4,00 & \checkmark \\
        \bottomrule
    \end{tabular}
\end{table}

\section{Identificação de Falhas de Aquisição}

A análise inicial de integridade identificou ocorrências de falhas em todos os 16 sujeitos do conjunto de dados, conforme consolidado na Tabela \ref{tab:problemas}. Tais anomalias são decorrentes de limitações instrumentais ou instabilidades na interface sensor-pele.

\begin{table}[htbp]
    \centering
    \caption{Resumo das falhas de integridade identificadas}
    \label{tab:problemas}
    \begin{tabular}{lc}
        \toprule
        \textbf{Categoria de Falha} & \textbf{Total de Ocorrências} \\
        \midrule
        Canais planos (ausência de sinal) & 32 \\
        Clipping (saturação de amplitude) & 66 \\
        Canais com alto nível de ruído & 16 \\
        \bottomrule
    \end{tabular}
    \end{table}

Conforme detalhado na Tabela \ref{tab:problemas}, identificaram-se 114 canais com algum tipo de falha de aquisição. O \textit{clipping} foi a anomalia mais frequente, representando 57,9\% das ocorrências totais, o que sugere que o elevado \textit{offset} DC do sistema Emotiv EPOC+ posiciona o sinal nas extremidades da faixa dinâmica do conversor A/D. Canais planos e ruidosos totalizaram 42,1\% das falhas, indicando problemas na interface sensor-pele e interferências ambientais.

\section{Inspeção Visual dos Registros Brutos}

A visualização dos sinais brutos para o sujeito 000, abrangendo as quatro modalidades coletadas, é ilustrada na Figura \ref{fig:sinais_brutos_s000}. Para uma avaliação comparativa da qualidade entre participantes, a Figura \ref{fig:visao_geral} apresenta os registros do canal AF7 de EEG para todos os 16 sujeitos do experimento.

\begin{figure}[htbp]
    \centering
    \includegraphics[width=0.95\textwidth]{../output/stage1_acquisition/figures/s000_raw_signals.png}
    \caption{Visualização dos biossinais brutos para o sujeito 000: canais de EEG, ECG, EMG e fNIRS.}
    \label{fig:sinais_brutos_s000}
\end{figure}

\begin{figure}[htbp]
    \centering
    \includegraphics[width=0.95\textwidth]{../output/stage1_acquisition/figures/overview_all_subjects.png}
    \caption{Panorama comparativo dos sinais de EEG (canal frontal AF7) para os 16 sujeitos do dataset.}
    \label{fig:visao_geral}
\end{figure}

\section{Avaliação de Qualidade do Sinal}

No segundo estágio do pipeline, executou-se a avaliação automática de qualidade (SQI) através da análise segmentada dos sinais. 

\subsection{Calibração de Limiares por Modalidade}
Os limiares de rejeição foram estabelecidos conforme descrito na 
Seção \ref{sec:calibracao_limiares}, calibrados para assegurar a robustez 
estatística das janelas de sinal. Observa-se que a modalidade fNIRS foi 
configurada com parâmetros de SNR extremamente permissivos devido à natureza 
do ruído intrínseco, enquanto o EEG exige maior estabilidade espectral.

A Tabela \ref{tab:sqi_thresholds} apresenta os limiares de rejeição 
estabelecidos para cada modalidade.

\begin{table}[htbp]
    \centering
    \caption{Limiares de rejeição por modalidade de biossinal}
    \label{tab:sqi_thresholds}
    \begin{tabular}{lcccc}
        \toprule
        \textbf{Modalidade} & \textbf{SNR mín. (dB)} & \textbf{Curtose máx.} & \textbf{Entropia mín.} \\
        \midrule
        EEG & $-5,0$ & $50,0$ & $0,15$ \\
        ECG & $-3,0$ & $100,0$ & $0,30$ \\
        EMG & $-5,0$ & $50,0$ & $0,20$ \\
        fNIRS & $-50,0$ & $100,0$ & $0,01$ \\
        \bottomrule
    \end{tabular}
\end{table}

\subsection{Análise Quantitativa de Rejeição}
A análise consolidada das taxas de rejeição por modalidade, considerando a totalidade dos 16 sujeitos, é apresentada na Tabela \ref{tab:sqi_summary}. Os resultados revelam disparidades significativas na estabilidade dos registros entre as modalidades. A taxa de rejeição do EEG ($\sim 5,0\%$) indica alta confiabilidade do sistema Emotiv EPOC+, enquanto o EMG apresentou perdas em mais de metade dos segmentos coletados ($\sim 56{,}0\%$). Por sua vez, o fNIRS alcançou taxa de rejeição de $97{,}5\%$, indicando que esta modalidade não possui integridade mínima para processamento.

\begin{table}[htbp]
    \centering
    \caption{Sumário das taxas de rejeição por modalidade}
    \label{tab:sqi_summary}
    \begin{tabular}{lccc}
        \toprule
        \textbf{Modalidade} & \textbf{Segmentos Rejeitados} & \textbf{Taxa de Rejeição} & \textbf{Status} \\
        \midrule
        EEG & 20/329 & $6{,}1\%$ & Estável \\
        ECG & 118/334 & $35{,}3\%$ & Crítico \\
        EMG & 187/334 & $56{,}0\%$ & Instável \\
        fNIRS & 306/314 & $97{,}5\%$ & Inutilizável \\
        \bottomrule
    \end{tabular}
\end{table}

\subsection{Inspeção de Artefatos Identificados}
A Figura \ref{fig:sqi_comparison} ilustra a diferenciação automática entre segmentos aceitos e rejeitados para a modalidade de EEG do sujeito 000. O algoritmo foi capaz de identificar com precisão o artefato de alta curtose causado por desvio abrupto na amplitude, preservando os segmentos com morfologia fisiológica característica.

\begin{figure}[htbp]
    \centering
    \includegraphics[width=0.95\textwidth]{../output/stage2_sqi/figures/sqi_comparison_000_eeg.png}
    \caption{Avaliação de qualidade (SQI) para EEG: comparação entre segmento aceito e segmento rejeitado por alta curtose (instrumental).}
    \label{fig:sqi_comparison}
\end{figure}

\subsection{Comparação com Abordagens da Literatura}
\label{sec:comparacao_literatura}

A Tabela \ref{tab:comparacao_sqi} apresenta a comparação entre a metodologia 
proposta e abordagens consolidadas da literatura para avaliação de qualidade 
de biossinais.

\begin{table}[htbp]
    \centering
    \caption{Comparação entre metodologias de SQI}
    \label{tab:comparacao_sqi}
    \begin{tabular}{lccc}
        \toprule
        \textbf{Trabalho} & \textbf{Modalidade} & \textbf{Métricas Utilizadas} & \textbf{Abordagem} \\
        \midrule
        \cite{cuena2019tool} & ECG & Curtose, Assimetria & hOS-SQI \\
        \cite{zheng2019eeg} & EEG & 6 métricas combinadas & EEGQI Framework \\
        \cite{taylor2022robustness} & ECG & SQIp, SQIsnr, SQIkur & Análise de Robustez \\
        \textbf{Proposto} & EEG/ECG/EMG & SNR, Curtose, Entropia & Thresholds Empíricos \\
        \bottomrule
    \end{tabular}
\end{table}

\section{Análise Estatística dos Sinais}
\label{sec:resultados_estatistica}

No terceiro estágio do pipeline, realizou-se a caracterização estatística abrangente dos registros biosignais. A Tabela \ref{tab:resumo_estatistico} apresenta os descritores consolidados para o primeiro sujeito (S000), demonstrando a variabilidade entre modalidades.

\begin{table}[htbp]
    \centering
    \caption{Caracterização estatística descritiva (Sujeito 000)}
    \label{tab:resumo_estatistico}
    \begin{tabular}{lccccc}
        \toprule
        \textbf{Modalidade} & \textbf{Canal} & \textbf{Média} & \textbf{Desvio Padrão} & \textbf{Assimetria} & \textbf{Curtose} \\
        \midrule
        EEG & AF7 & $-64254{,}6$ & $3266{,}8$ & $1{,}89$ & $4{,}59$ \\
        EEG & AF8 & $-67520{,}9$ & $3297{,}9$ & $1{,}89$ & $4{,}91$ \\
        ECG & ECG & $-530{,}3$ & $799{,}9$ & $200{,}2$ & $47055{,}6$ \\
        EMG & EMG & $8530{,}3$ & $227{,}5$ & $-0{,}26$ & $30{,}23$ \\
        fNIRS & HbO & $4{,}1 \times 10^{-9}$ & $3{,}5 \times 10^{-7}$ & $-9{,}01$ & $102{,}87$ \\
        \bottomrule
    \end{tabular}
\end{table}

\subsection{Resultados dos Testes de Normalidade}
O teste de Shapiro-Wilk revelou que 100\% (192/192) dos canais analisados em todos os 16 sujeitos divergem significativamente de uma distribuição normal ($p < 0,05$). Tais achados são corroborados pelos altos valores de curtose e assimetria, especialmente na modalidade ECG (Curtose $\approx 47055{,}6$), evidenciando a presença de picos R proeminentes e ruído impulsivo. A Figura \ref{fig:histograma_eeg} ilustra a distribuição de amplitude para o canal AF7 de EEG, comparando-a à curva normal teórica.

\begin{figure}[htbp]
    \centering
    \includegraphics[width=0.85\textwidth]{../output/stage3_statistics/figures/stat_histogram_000_eeg_AF7.png}
    \caption{Histograma de amplitudes do canal EEG-AF7 (Sujeito 000) com sobreposição de curva normal teórica.}
    \label{fig:histograma_eeg}
\end{figure}

Para avaliar visualmente o desvio da normalidade e a dispersão entre canais, a Figura \ref{fig:qq_boxplot} apresenta o gráfico Q-Q (\textit{Quantile-Quantile}) e o \textit{boxplot} comparativo para a modalidade EEG.

\begin{figure}[htbp]
    \centering
    \begin{subfigure}[b]{0.48\textwidth}
        \centering
        \includegraphics[width=\textwidth]{../output/stage3_statistics/figures/stat_qq_000_eeg_AF7.png}
        \caption{Gráfico Q-Q (AF7)}
        \label{fig:qq_eeg}
    \end{subfigure}
    \hfill
    \begin{subfigure}[b]{0.48\textwidth}
        \centering
        \includegraphics[width=\textwidth]{../output/stage3_statistics/figures/stat_boxplot_000_eeg.png}
        \caption{Boxplot por Canal}
        \label{fig:boxplot_eeg}
    \end{subfigure}
    \caption{Visualizações de diagnóstico estatístico para EEG (Sujeito 000): (a) avaliação de normalidade via quantis e (b) dispersão e \textit{outliers} entre os 8 canais.}
    \label{fig:qq_boxplot}
\end{figure}

\subsection{Testes de Homocedasticidade}
A análise de homocedasticidade pelo teste de Levene entre os 8 canais de EEG demonstrou diferenças significativas de variância ($W = 7048{,}7, p < 0,05$). Este resultado sugere que os sensores de EEG, apesar de serem da mesma tecnologia (Emotiv EPOC+), apresentam níveis de ruído instrumental e impedância sensor-pele heterogêneos entre os locais de escalpo. Para a modalidade fNIRS, a falta de homocedasticidade entre HbO e HbR ($p \approx 3{,}1 \times 10^{-39}$) reforça a instabilidade desta aquisição.

\subsection{Correlação entre Canais e Modalidades}
A matriz de correlação cruzada, apresentada na Figura \ref{fig:heatmap_correlacao}, revela fortes interdependências entre os canais de EEG frontais (AF7 e AF8, $r > 0,80$), conforme esperado pela proximidade física dos sensores. Observou-se correlação desprezível entre as modalidades (ex: EEG vs ECG), indicando que os artefatos de pulsação cardíaca no registro de EEG não são dominantes no estado bruto.

\begin{figure}[htbp]
    \centering
    \includegraphics[width=0.9\textwidth]{../output/stage3_statistics/figures/stat_correlation_000.png}
    \caption{Matriz de correlação cruzada entre canais e modalidades para o Sujeito 000.}
    \label{fig:heatmap_correlacao}
\end{figure}

Para avaliar a estabilidade das correlações em nível populacional, a Figura \ref{fig:heatmap_global} apresenta a matriz agregada considerando o conjunto total de sujeitos analisados.

\begin{figure}[htbp]
    \centering
    \includegraphics[width=0.9\textwidth]{../output/stage3_statistics/figures/stat_correlation_999.png}
    \caption{Matriz de correlação global agregada para todos os sujeitos do dataset.}
    \label{fig:heatmap_global}
\end{figure}


\section{Limpeza e Correção de Sinais}
\label{sec:resultados_limpeza}

No quarto estágio do pipeline, aplicou-se o conjunto de técnicas de filtragem e normalização conforme descrito na Seção \ref{sec:metodo_limpeza}. A Tabela \ref{tab:cleaning_summary} apresenta o resumo consolidado dos efeitos da limpeza para o sujeito 000, evidenciando reduções drásticas de variância e normalização de curtose.

\begin{table}[htbp]
    \centering
    \caption{Efeitos da limpeza de sinais (Sujeito 000)}
    \label{tab:cleaning_summary}
    \resizebox{\textwidth}{!}{
        \begin{tabular}{lccccc}
            \toprule
            \textbf{\textit{Modalidade}} & \textbf{Var. Redução (\%)} & \textbf{Curtose (antes)} & \textbf{Curtose (depois)} & \textbf{SNR (dB)} & \textbf{Amostras Winsorizadas} \\
            \midrule
            EEG (AF7) & 99,9 & 4,60 & 3,44 & 12,36 & 7.418 \\
            EEG (PO7) & 99,9 & 221,83 & -1,12 & -2,06 & 7.418 \\
            ECG & 94,0 & 47055,61 & 1,30 & 6,94 & 6.488 \\
            EMG & 99,8 & 30,23 & -0,95 & -1,26 & 6.488 \\
            fNIRS (HbO) & 97,2 & 102,87 & 1,61 & -60,00 & 232 \\
            \bottomrule
        \end{tabular}
    }
\end{table}

\subsection{Redução de Variância e Efeito de Cohen}

Observou-se que os coeficientes $d$ de Cohen calculados entre os estados antes e depois da limpeza indicam efeito de tamanho grande ($|d| > 0{,}8$) para todas as modalidades principais (EEG, ECG, EMG), confirmando que o processo de limpeza altera substancialmente a distribuição dos sinais. A variância reduziu-se em mais de 99\% para EEG e EMG, e 94\% para ECG, evidenciando a eficácia da combinação de filtragem e Winsorização na atenuação de ruído instrumental.

\subsection{Normalização de Curtose}

A curtose, que originalmente apresentava valores extremos (especialmente em ECG com $K \approx 47000$ devido aos picos R), foi normalizada para valores próximos de zero após a limpeza ($|K| < 3$), indicando distribuições aproximadamente mesocúrticas. Para o canal PO7 de EEG, a curtose passou de 221,83 para -1,12, evidenciando a remoção de \textit{spikes} instrumentais.

\subsection{Interpolação de Lacunas}

Identificou-se que os sinais não apresentaram lacunas significativas (gaps $> 1$s) que requeressem interpolação nos dados do sujeito 000. A função de interpolação foi executada corretamente em todos os demais sujeitos.

\subsection{Manutenção de fNIRS}

Conforme planejado na discussão do Estágio 3, a modalidade fNIRS foi mantida no pipeline de limpeza apesar de apresentar SNR constante de -60 dB e efeito de tamanho negligível. A justificativa para manutenção é que a filtragem passa-banda (0,01--0,1 Hz) reduz a variância em mais de 97\% (Tabela \ref{tab:cleaning_summary}), sugerindo que mesmo um sinal ruidoso pode beneficiar-se do pré-processamento para análises de tendências lentas.

Esta análise foi conduzida como prova de conceito metodológico, demonstrando a eficácia dos filtros mesmo em sinais de baixa qualidade. No entanto, para as fases subsequentes de extração de atributos e classificação (Estágios 6--10), o fNIRS será excluído do pipeline devido à sua taxa de rejeição superior a 95\%.

\section{Segmentação Temporal dos Sinais}
\label{sec:resultados_segmentacao}

No quinto estágio do pipeline, executou-se a segmentação temporal dos sinais limpos em janelas de 5 segundos com 50\% de sobreposição. A Tabela \ref{tab:segmentacao_summary} apresenta o consolidado da segmentação por modalidade.

\begin{table}[htbp]
    \centering
    \caption{Sumário da segmentação temporal por modalidade}
    \label{tab:segmentacao_summary}
    \begin{tabular}{lccc}
        \toprule
        \textbf{Modalidade} & \textbf{Janelas Totais} & \textbf{Janelas Válidas} & \textbf{Taxa de Retenção} \\
        \midrule
        EEG (8 canais) & 329 & 309 & $93{,}9\%$ \\
        ECG & 334 & 216 & $64{,}7\%$ \\
        EMG & 334 & 147 & $44{,}0\%$ \\
        fNIRS & 314 & 8 & $2{,}5\%$ \\
        \bottomrule
    \end{tabular}
\end{table}

\subsection{Análise de Estabilidade Intra-Janela}

O teste de estacionariedade ADF e o coeficiente de variação (CV) foram aplicados a cada segmento para validar sua qualidade fisiológica. A Figura \ref{fig:window_stability} ilustra a distribuição dos valores de CV por modalidade para o EEG.

\begin{figure}[htbp]
    \centering
    \includegraphics[width=0.9\textwidth]{../output/stage5_segmentation/figures/window_stability_000_eeg.png}
    \caption{Distribuição do coeficiente de variação (CV) para EEG (Sujeito 000).}
    \label{fig:window_stability}
\end{figure}

A Tabela \ref{tab:segmentacao_metrics} apresenta as métricas de estabilidade consolidadas para o sujeito 000, demonstrando que o EEG apresenta excelente uniformidade entre janelas (CV médio de 0,41\%), enquanto o ECG apresenta maior variabilidade de amplitude (CV médio de 45,9\%). A estacionariedade ADF indica que o ECG é fortemente estacionário (100\%) devido à periodicidade cardíaca, enquanto o EEG é predominantemente não-estacionário (0\%), coerente com a natureza dinâmica dos potenciais corticais.

\begin{table}[htbp]
    \centering
    \caption{Métricas de estabilidade intra-janela (Sujeito 000)}
    \label{tab:segmentacao_metrics}
    \begin{tabular}{lcccc}
        \toprule
        \textbf{Modalidade} & \textbf{CV Médio (\%)} & \textbf{CV Máx (\%)} & \textbf{ADF Estacionário (\%)} \\
        \midrule
        EEG & $0{,}41$ & $3{,}21$ & $0{,}0$ \\
        ECG & $45{,}91$ & $73{,}45$ & $100{,}0$ \\
        EMG & $0{,}29$ & $0{,}53$ & $37{,}5$ \\
        fNIRS & $2045{,}44$ & $6067{,}80$ & $100{,}0$ \\
        \bottomrule
    \end{tabular}
\end{table}

\subsection{Variância Inter-Janelas}

A análise de variância entre segmentos consecutivos revelou que o EEG apresenta menor variabilidade inter-janelas (indicando segmentos temporalmente consistentes), enquanto o EMG apresenta maior dispersão, refletindo a natureza descontínua da atividade muscular. A Figura \ref{fig:inter_window} apresenta a análise comparativa para EEG.

\begin{figure}[htbp]
    \centering
    \includegraphics[width=0.9\textwidth]{../output/stage5_segmentation/figures/inter_window_variance_000_eeg.png}
    \caption{Variância inter-janelas para EEG (Sujeito 000).}
    \label{fig:inter_window}
\end{figure}

\subsection{Propagação de Máscaras de Qualidade}

As máscaras de rejeição geradas no Estágio 2 foram propagadas para o janelamento, resultando nas taxas de rejeição apresentadas na Tabela \ref{tab:segmentacao_summary}. Observou-se que:

\begin{itemize}
    \item \textbf{EEG:} 20 janelas rejeitadas (6,1\%), principalmente por artefatos de piscada ocular e artefatos de movimento.
    \item \textbf{ECG:} 118 janelas rejeitadas (35,3\%), principalmente por baixa SNR e variabilidade excessiva.
    \item \textbf{EMG:} 187 janelas rejeitadas (56,0\%), indicando instabilidade na interface sensor-pele durante o protocolo de indução emocional.
    \item \textbf{fNIRS:} 306 janelas rejeitadas (97,5\%), confirmando a exclusão desta modalidade do pipeline.
\end{itemize}

\subsection{Visualização das Janelas}

A Figura \ref{fig:segmentation_windows} apresenta a representação visual das janelas segmentadas para o sujeito 000 (EEG), demonstrando a aplicação correta do janelamento com sobreposição.

\begin{figure}[htbp]
    \centering
    \includegraphics[width=0.9\textwidth]{../output/stage5_segmentation/figures/segmentation_windows_000_eeg.png}
    \caption{Segmentação temporal com janelas de 5s e 50\% de sobreposição para EEG (Sujeito 000).}
    \label{fig:segmentation_windows}
\end{figure}

\subsection{Cache de Segmentos Gerados}

Ao final do estágio, foram gerados 64 arquivos \texttt{.npz} (4 modalidades × 16 sujeitos) contendo os segmentos válidos. Cada arquivo preserva:
\begin{itemize}
    \item Sinais particionados com metadados temporais.
    \item Máscaras de qualidade propagadas do SQI.
    \item Métricas de estabilidade por janela.
\end{itemize}

Estes arquivos constituem a entrada para o Estágio 6 de extração de atributos, otimizando o tempo de processamento ao evitar reprocessamento de sinais brutos.

\subsection{Análise Inter-Sujeitos e Padrões Observados}

A análise da variabilidade entre os 16 sujeitos revelou padrões distintos por modalidade, conforme apresentado na Tabela \ref{tab:subjects_segmentacao}. As observações mais importantes revelam também particularidades de cada sujeito para com os sinais coletados, além de análise inter-modalidades e inter-janelas.

\begin{table}[htbp]
    \centering
    \caption{Retenção de janelas por sujeito e modalidade}
    \label{tab:subjects_segmentacao}
    \resizebox{\textwidth}{!}{
        \begin{tabular}{
            l
            S[table-format=2.0]
            S[table-format=2.0, table-space-text-post={\ (100\%)}]
            @{\hspace{4ex}}
            S[table-format=2.0]
            S[table-format=2.0, table-space-text-post={\ (100\%)}]
            @{\hspace{4ex}}
            S[table-format=2.0]
            S[table-format=2.0, table-space-text-post={\ (100\%)}]
            @{\hspace{4ex}}
            S[table-format=2.0]
            S[table-format=2.0, table-space-text-post={\ (100\%)}]
        }
            \toprule
            \textbf{Sujeito}
                & \multicolumn{2}{c}{\textbf{EEG}}
                & \multicolumn{2}{c}{\textbf{ECG}}
                & \multicolumn{2}{c}{\textbf{EMG}}
                & \multicolumn{2}{c}{\textbf{fNIRS}} \\
            \cmidrule(lr){2-3}\cmidrule(lr){4-5}\cmidrule(lr){6-7}\cmidrule(lr){8-9}
            & {\textbf{Total}} & {\textbf{Válidas}}
            & {\textbf{Total}} & {\textbf{Válidas}}
            & {\textbf{Total}} & {\textbf{Válidas}}
            & {\textbf{Total}} & {\textbf{Válidas}} \\
            \midrule
            s000 & 28 & 23 { (82\%)} & 51 & 16 { (31\%)} & 51 & 24 { (47\%)} & 28 & 4 { (14\%)} \\
            s001 & 23 & 22 { (96\%)} & 20 & 17 { (85\%)} & 20 &  7 { (35\%)} & 23 & 2 { (9\%)}  \\
            s002 & 15 & 14 { (93\%)} & 15 & 14 { (93\%)} & 15 &  0 { (0\%)}  & 15 & 0 { (0\%)}  \\
            s003 & 16 & 15 { (94\%)} & 16 &  2 { (12\%)} & 16 & 13 { (81\%)} & 16 & 0 { (0\%)}  \\
            s004 & 20 & 19 { (95\%)} & 16 &  4 { (25\%)} & 16 & 10 { (62\%)} & 20 & 0 { (0\%)}  \\
            s005 & 21 & 20 { (95\%)} & 18 & 17 { (94\%)} & 18 & 17 { (94\%)} & 21 & 0 { (0\%)}  \\
            s006 & 20 & 19 { (95\%)} & 19 &  2 { (11\%)} & 19 &  8 { (42\%)} & 20 & 1 { (5\%)}  \\
            s007 & 16 & 15 { (94\%)} & 17 &  9 { (53\%)} & 17 &  2 { (12\%)} & 16 & 0 { (0\%)}  \\
            s008 & 19 & 18 { (95\%)} & 20 & 19 { (95\%)} & 20 &  8 { (40\%)} & 19 & 0 { (0\%)}  \\
            s009 & 20 & 19 { (95\%)} & 20 & 19 { (95\%)} & 20 &  3 { (15\%)} & 20 & 0 { (0\%)}  \\
            s010 & 14 & 13 { (93\%)} & 15 & 14 { (93\%)} & 15 &  4 { (27\%)} & 14 & 0 { (0\%)}  \\
            s011 & 24 & 23 { (96\%)} & 24 & 23 { (96\%)} & 24 & 21 { (88\%)} & 24 & 0 { (0\%)}  \\
            s012 & 23 & 22 { (96\%)} & 24 & 23 { (96\%)} & 24 & 10 { (42\%)} & 23 & 1 { (4\%)}  \\
            s013 & 26 & 25 { (96\%)} & 20 &  0 { (0\%)}  & 20 & 10 { (50\%)} & 26 & 0 { (0\%)}  \\
            s014 & 25 & 24 { (96\%)} & 23 & 22 { (96\%)} & 23 & 10 { (43\%)} & 25 & 0 { (0\%)}  \\
            s015 & 19 & 18 { (95\%)} & 16 & 15 { (94\%)} & 16 &  0 { (0\%)}  &  4 & 0 { (0\%)}  \\
            \bottomrule
        \end{tabular}
    }
\end{table}

A análise por sujeito revela padrões de qualidade heterogêneos entre participantes. O sujeito \textbf{s005} apresentou o melhor perfil global (EEG 95\%, ECG 94\%, EMG 94\%), enquanto \textbf{s011} foi o único a superar 88\% em todas as três modalidades ativas. Em contraste, \textbf{s000} registrou as maiores perdas multimodais simultâneas: EEG 82\% (menor entre todos), ECG 31\% e EMG 47\%, sugerindo problemas de fixação do equipamento durante o protocolo.

No ECG, três sujeitos apresentaram rejeição quase total: \textbf{s003} (88\% rejeitado), \textbf{s006} (89\%) e, criticamente, \textbf{s013} (100\% rejeitado), que foi excluído do dataset cruzado final por ausência de janelas ECG válidas. Esses casos indicam provável desconexão do eletrodo ou interferência de movimento de alta amplitude. No EMG, os sujeitos \textbf{s002} e \textbf{s015} tiveram 100\% de rejeição — falhas completas de aquisição que os excluem igualmente do dataset cruzado. Os sujeitos \textbf{s007} (88\% rejeitado) e \textbf{s009} (85\% rejeitado) apresentaram EMG com baixíssima retenção, consistente com a ausência de tarefa motora explícita no protocolo. Para o fNIRS, todos os 16 sujeitos tiveram retenção inferior a 15\%, confirmando a exclusão definitiva desta modalidade.


\begin{enumerate}
    \item \textbf{EEG:} A taxa de retenção foi consistentemente alta (82--96\%), com média de 93,9\%. O sujeito s000 apresentou a menor taxa (82\%), possivelmente devido a artefatos de movimento prolongados durante o protocolo. Os coeficientes de variação (CV) foram extremamente baixos (média de 0,41\% para s000), indicando estabilidade temporal excepcional do sinal EEG.

    \item \textbf{ECG:} Observou-se alta variabilidade inter-sujeitos. Quatro sujeitos (s003, s004, s006, s013) apresentaram taxa de rejeição superior a 70\%, sugerindo problemas crônicos de contato sensor ou arritmias cardíacas. Em contraste, seis sujeitos (s001, s002, s005, s008, s009, s011) alcançaram retenção superior a 85\%. O teste ADF indicou que 97,7\% das janelas de ECG são estacionárias, coerente com a periodicidade cardíaca quase-determinística.

    \item \textbf{EMG:} Dois sujeitos (s002, s015) apresentaram 0\% de janelas utilizáveis, indicando falhas completas na aquisição. A taxa média de estacionariedade ADF foi baixa (14,3\%), refletindo a natureza descontínua e não-estacionária da atividade eletromiográfica durante a indução emocional.

    \item \textbf{fNIRS:} Confirmou-se a exclusão de fNIRS para todos os sujeitos, com apenas 8 janelas válidas (2,5\%) no dataset completo. Os valores de CV foram extremos (média de 2045\%), indicando variabilidade não-fisiológica nos dados.

    \item \textbf{Correlação Inter-Modalidades:} A correlação de Pearson entre as taxas de retenção de diferentes modalidades foi fraca (EEG-ECG: $r = 0{,}24$; ECG-EMG: $r = -0{,}25$), indicando que sujeitos com boa qualidade de EEG não necessariamente apresentam boa qualidade de ECG ou EMG. Isto sugere que as fontes de artefato são específicas por modalidade e não correlacionadas com o estado geral do participante.

    \item \textbf{Variância Inter-Janelas:} O ECG apresentou maior variância entre janelas consecutivas ($\sigma^2 = 0{,}012$) em comparação com EEG ($\sigma^2 = 0{,}003$) e EMG ($\sigma^2 = 0{,}0004$), refletindo a maior dinâmica temporal do sinal cardíaco.
\end{enumerate}

\section{Discussão dos Achados}

A análise integrada dos resultados indica que os sinais de EEG apresentam a maior estabilidade qualitativa, com baixas taxas de rejeição. No entanto, as modalidades de ECG e EMG demonstraram alta dependência da qualidade do contato sensor-pele e de artefatos de movimento, resultando em perdas significativas de janelas temporais. 

A rejeição de hipótese nula nos testes de normalidade (Shapiro-Wilk) e homocedasticidade (Levene) justifica a adoção de estratégias de processamento robusto nos estágios subsequentes. Em particular, a filtragem espacial e a normalização de variância (Z-score por canal) tornam-se indispensáveis antes da extração de atributos, para mitigar os vieses introduzidos pela heterogeneidade instrumental e pelos artefatos de cauda pesada.

A modalidade de fNIRS apresentou taxa de rejeição próxima à totalidade ($>95\%$), confirmando falhas estruturais na aquisição original do dataset (distâncias entre fonte e detector nulas). Tais achados, aliados à curtose anômala ($>100$), justificam a exclusão definitiva do fNIRS nas etapas de filtragem (Estágio 4) e extração de atributos (Estágio 6), otimizando o pipeline para os biossinais que mantêm integridade fisiológica mínima.

\section{Extração de Atributos}
\label{sec:resultados_extracao}

O sexto estágio do pipeline executou a extração de atributos sobre os 64 arquivos de segmentos gerados no Estágio 5, processando as três modalidades ativas (EEG, ECG e EMG) para todos os 16 sujeitos.

\subsection{Matriz de Atributos Gerada}

Ao total, foram extraídas 2.780 observações (janela $\times$ canal), distribuídas conforme a Tabela \ref{tab:features_summary}.

\begin{table}[htbp]
    \centering
    \caption{Resumo da matriz de atributos extraída no Estágio 6}
    \label{tab:features_summary}
    \begin{tabular}{lccc}
        \toprule
        \textbf{Modalidade} & \textbf{Atributos} & \textbf{Observações} & \textbf{Descrição adicional} \\
        \midrule
        EEG & 18 & 2.417 & 9 temporais + 4 espectrais + 5 bandas \\
        ECG & 20 & 216  & 13 gerais + 7 HRV (tempo e freq.) \\
        EMG & 13 & 147  & 9 temporais + 4 espectrais \\
        \midrule
        \textbf{Total} & --- & \textbf{2.780} & 3 modalidades, 16 sujeitos \\
        \bottomrule
    \end{tabular}
\end{table}

\subsection{Atributos Temporais e Espectrais}

A Figura \ref{fig:feature_dist_eeg} apresenta os boxplots dos atributos de EEG agregados sobre todos os sujeitos e janelas. Observa-se elevada dispersão nos atributos de potência de banda, refletindo a variabilidade inter-sujeito característica dos registros de EEG emocional.

\begin{figure}[htbp]
    \centering
    \includegraphics[width=0.95\textwidth]{../output/stage6_features/figures/feature_distributions_eeg.png}
    \caption{Distribuição dos atributos de EEG (boxplots) agregados sobre todos os sujeitos e janelas válidas.}
    \label{fig:feature_dist_eeg}
\end{figure}

\begin{figure}[htbp]
    \centering
    \includegraphics[width=0.95\textwidth]{../output/stage6_features/figures/feature_distributions_ecg.png}
    \caption{Distribuição dos atributos de ECG, incluindo os atributos de HRV no domínio do tempo e da frequência.}
    \label{fig:feature_dist_ecg}
\end{figure}

\begin{figure}[htbp]
    \centering
    \includegraphics[width=0.95\textwidth]{../output/stage6_features/figures/feature_distributions_emg.png}
    \caption{Distribuição dos atributos de EMG, com ênfase nos descritores temporais e espectrais.}
    \label{fig:feature_dist_emg}
\end{figure}

\subsection{Correlação entre Atributos}

A Figura \ref{fig:feature_corr_eeg} exibe o mapa de calor das correlações de Pearson entre os atributos de EEG. Destaca-se a alta correlação entre \textit{hjorth\_activity} e variância ($r \approx 1{,}00$), o que é esperado pela definição dos parâmetros de Hjorth, e entre as potências das bandas adjacentes, evidenciando a continuidade espectral dos sinais.

\begin{figure}[htbp]
    \centering
    \includegraphics[width=0.90\textwidth]{../output/stage6_features/figures/feature_correlation_eeg.png}
    \caption{Mapa de calor das correlações de Pearson entre atributos de EEG.}
    \label{fig:feature_corr_eeg}
\end{figure}

\subsection{Potência por Bandas de EEG}

A Figura \ref{fig:eeg_bands} apresenta as potências médias por banda de frequência para cada sujeito. A banda beta (13--30 Hz) dominou com potência média de 41,7 $\mu V^2$/Hz, seguida pela banda alpha (8--13 Hz) com 19,0 $\mu V^2$/Hz. Esta distribuição é consistente com a literatura de EEG emocional, onde a banda beta é associada ao processamento cognitivo e à resposta ao estímulo visual. Destaca-se que o sujeito s000 apresenta potência na banda delta aproximadamente 80 vezes superior à mediana dos demais sujeitos, indicando a presença de artefatos de baixa frequência persistentes neste participante — artefato já identificado no Estágio 1 como canal com alto nível de ruído.

\begin{figure}[htbp]
    \centering
    \includegraphics[width=0.95\textwidth]{../output/stage6_features/figures/eeg_band_power_summary.png}
    \caption{Potência média por banda de frequência (delta, theta, alpha, beta, gamma) para cada sujeito.}
    \label{fig:eeg_bands}
\end{figure}

\subsection{Variabilidade da Frequência Cardíaca (ECG)}

Os atributos de HRV extraídos para o ECG revelaram padrões fisiologicamente plausíveis. O intervalo RR médio foi de $731{,}3 \pm 94{,}8$ ms (correspondente a $\approx 82$ bpm), onde o desvio padrão reflete a variação entre janelas; o SDNN médio intra-janela foi de $105{,}4$ ms e o RMSSD médio de $132{,}4$ ms. A alta variabilidade observada no RMSSD é consistente com a natureza dos registros de curta duração (janelas de 5 s), onde a estimativa HRV tem maior incerteza estatística. Os atributos espectrais de HRV (\textit{lf\_power}, \textit{hf\_power} e \textit{lf\_hf\_ratio}) não foram computáveis neste estágio: janelas de 5 s fornecem resolução espectral de apenas 0{,}2 Hz, insuficiente para resolver as bandas LF (0{,}04--0{,}15 Hz) e HF (0{,}15--0{,}4 Hz) da análise espectral de HRV, que requer registros de pelo menos 2 minutos \cite{taskforce1996}.

\subsection{Atributos de EMG}

Para o EMG, o RMS médio foi de 5.971 $\mu V$, a frequência média de 27{,}4 Hz e a frequência mediana de 20{,}2 Hz. A baixa frequência dominante indica que a atividade muscular capturada é principalmente de baixa intensidade, coerente com o protocolo experimental de indução emocional passiva (sem tarefa motora explícita).

\section{Engenharia de Atributos}
\label{sec:resultados_engenharia}

O sétimo estágio enriqueceu a matriz de atributos com novos descritores derivados, normalizações e análises de qualidade, conforme descrito na Seção \ref{sec:metodo_engenharia}.

\subsection{Matriz Enriquecida}

A Tabela \ref{tab:eng_summary} resume o número de colunas adicionadas por modalidade após a engenharia de atributos.

\begin{table}[htbp]
    \centering
    \caption{Colunas adicionadas pela engenharia de atributos por modalidade}
    \label{tab:eng_summary}
    \begin{tabular}{lcccc}
        \toprule
        \textbf{Modalidade} & \textbf{Orig.} & \textbf{+ Razões} & \textbf{+ Norm.} & \textbf{+ $\Delta$/$\Delta^2$} \\
        \midrule
        EEG & 18 & 5 & 23 & 10 \\
        ECG & 20 & — & 20 & 10 \\
        EMG & 13 & — & 13 & 10 \\
        \bottomrule
    \end{tabular}
\end{table}

\subsection{Razões de Potência de Banda EEG}

A Figura \ref{fig:band_ratios} apresenta a distribuição das razões de potência por fase experimental. O índice alpha/beta foi o atributo de maior poder discriminante entre fases ($F = 23{,}16$, $p < 0{,}001$), seguido pela razão theta/beta ($F = 17{,}41$, $p < 0{,}001$). Durante a fase de estimulação observa-se redução do índice alpha/beta em relação à linha de base, coerente com o aumento da atividade cortical reportado na literatura de EEG emocional.

\begin{figure}[htbp]
    \centering
    \includegraphics[width=0.95\textwidth]{../output/stage7_engineering/figures/band_ratios_eeg.png}
    \caption{Razões de potência de banda EEG por fase experimental.}
    \label{fig:band_ratios}
\end{figure}

\subsection{Discriminabilidade entre Fases}

A Figura \ref{fig:discriminability} apresenta as $F$-estatísticas ANOVA para os 20 atributos mais discriminantes de cada modalidade. Para o EEG, 6 atributos foram estatisticamente significativos ($p < 0{,}05$): \textit{alpha\_beta\_ratio} ($F = 23{,}16$), \textit{theta\_beta\_ratio} ($F = 17{,}41$), \textit{mav} ($F = 9{,}32$), \textit{rms} ($F = 9{,}32$), \textit{mean\_freq} ($F = 7{,}78$) e \textit{spectral\_entropy} ($F = 4{,}78$), sendo os dois primeiros as razões espectrais recém-criadas. Para ECG e EMG, nenhum atributo atingiu significância estatística ($F < 3{,}0$), o que é consistente com o protocolo de indução emocional passiva, que não induz variações autonômicas ou musculares sistemáticas suficientes para discriminar fases em janelas de 5 s.

\begin{figure}[htbp]
    \centering
    \includegraphics[width=0.95\textwidth]{../output/stage7_engineering/figures/phase_discriminability.png}
    \caption{Discriminabilidade de fase (ANOVA $F$-estatística) por atributo e modalidade.}
    \label{fig:discriminability}
\end{figure}

\subsection{Análise de Redundância}

A análise de correlação de Pearson identificou 7 pares redundantes no EEG ($|r| \geq 0{,}95$), 2 no EMG e nenhum no ECG, conforme ilustrado na Figura \ref{fig:redundancy_eeg}. Os pares redundantes no EEG concentram-se entre \textit{hjorth\_activity} e variância (correlação perfeita por definição) e entre as potências de bandas espectralmente adjacentes. Estes pares serão candidatos à remoção no Estágio 9 de seleção de atributos.

\begin{figure}[htbp]
    \centering
    \includegraphics[width=0.85\textwidth]{../output/stage7_engineering/figures/feature_redundancy_eeg.png}
    \caption{Mapa de calor de redundância entre atributos de EEG (correlação de Pearson).}
    \label{fig:redundancy_eeg}
\end{figure}

\section{Redução de Dimensionalidade}
\label{sec:resultados_reducao}

O oitavo estágio aplicou PCA sobre as matrizes agregadas do Estágio 7, conforme descrito na Seção~\ref{sec:metodo_reducao}. A Tabela~\ref{tab:pca_summary} resume os resultados por modalidade.

\begin{table}[htbp]
    \centering
    \caption{Resultados da PCA por modalidade}
    \label{tab:pca_summary}
    \begin{tabular}{cccccccc}
        \toprule
        \textbf{Modalidade} & \textbf{Obs.} & \makecell[c]{\textbf{Atrib.} \\ \textbf{entrada}} & \textbf{Descartados} & \textbf{PC@90\%} & \textbf{PC@95\%} & \makecell[c]{\textbf{Var.} \\ \textbf{PC1}} & \makecell[c]{\textbf{Var.} \\ \textbf{PC2}} \\
        \midrule
        EEG & 375 & 220 & 4  & 26 & 38 & 18{,}7\% & 15{,}7\% \\
        ECG & 41  & 172 & 28 & 18 & 22 & 17{,}9\% & 11{,}3\% \\
        EMG & 37  & 140 & 4  & 15 & 18 & 18{,}9\% & 16{,}1\% \\
        \bottomrule
    \end{tabular}
\end{table}

\subsection{EEG}

Para o EEG, 375 observações (sujeito $\times$ canal $\times$ fase) foram decompostas a partir de 220 atributos (4 colunas \textit{zcr\_norm} descartadas por excesso de NaN). O primeiro componente explica 18{,}7\% da variância total e o segundo 15{,}7\%; os cinco primeiros componentes acumulam 54{,}0\% e os dez primeiros 71{,}1\%. São necessários 26 componentes para atingir 90\% e 38 para 95\% de variância explicada, representando uma compressão de 83\% (220 → 38 dimensões).

A análise das cargas do PC1 revela que os atributos de maior peso são derivados de segunda ordem dos poderes espectrais (\textit{delta2\_alpha\_power\_std}, \textit{delta2\_alpha\_power\_max}, \textit{delta2\_beta\_power\_std}), indicando que a principal fonte de variância entre observações é a instabilidade temporal dos poderes nas bandas alfa e beta. O PC2 é dominado pelos agregados normalizados de máximo (\textit{total\_power\_norm\_max}, \textit{theta\_power\_norm\_max}), capturando contrastes de amplitude entre sujeitos. A projeção das observações no plano PC1~$\times$~PC2 (Figura~\ref{fig:pca_scatter_eeg}) revela sobreposição parcial entre as três fases, consistente com os resultados de discriminabilidade do Estágio~7 e com a natureza distribuída da informação de fase no espaço PCA.

\begin{figure}[htbp]
    \centering
    \includegraphics[width=0.95\textwidth]{../output/stage8_dimreduction/figures/scree_plot_eeg.png}
    \caption{Diagrama de cotovelo para o EEG.}
    \label{fig:scree_eeg}
\end{figure}

\begin{figure}[htbp]
    \centering
    \includegraphics[width=0.95\textwidth]{../output/stage8_dimreduction/figures/cumulative_variance_eeg.png}
    \caption{Curva de variância acumulada para o EEG.}
    \label{fig:cumvar_eeg}
\end{figure}

\begin{figure}[htbp]
    \centering
    \includegraphics[width=0.80\textwidth]{../output/stage8_dimreduction/figures/pca_scatter_eeg.png}
    \caption{Dispersão PC1 $\times$ PC2 para EEG por fase experimental.}
    \label{fig:pca_scatter_eeg}
\end{figure}

\subsection{ECG e EMG}

Para o ECG, 28 das 200 colunas de entrada foram descartadas por excesso de NaN (24 relativas às bandas LF/HF e seus derivados, cujos valores são nulos devido à limitação de resolução espectral das janelas de 5~s). A decomposição resultou em 41 componentes (limitada pelo número de observações), com 18 componentes para 90\% e 22 para 95\% de variância. O PC1 do ECG é dominado por atributos de mobilidade de Hjorth e seus deltas (\textit{hjorth\_mobility\_std}, $r = 0{,}167$), refletindo variabilidade na morfologia do sinal cardíaco.

Para o EMG, 37 observações e 140 atributos (4 descartados) resultam em 15 e 18 componentes para os limiares de 90\% e 95\%, respectivamente. O EMG apresenta a maior variância nos primeiros cinco componentes (61{,}6\%), sugerindo menor redundância na estrutura dos atributos em relação ao EEG.

\begin{figure}[htbp]
    \centering
    \includegraphics[width=\textwidth]{../output/stage8_dimreduction/figures/pca_loadings_eeg.png}
    \caption{Mapa de cargas PCA para EEG: primeiros 5 componentes e os 15 atributos de maior carga absoluta.}
    \label{fig:pca_loadings_eeg}
\end{figure}

\section{Seleção de Atributos}
\label{sec:resultados_selecao}

O nono estágio aplicou três famílias de seleção sobre os componentes principais produzidos pelo Estágio 8, conforme descrito na Seção~\ref{sec:metodo_selecao}. A Tabela~\ref{tab:selection_summary} apresenta os resultados consolidados por modalidade.

\begin{table}[htbp]
    \centering
    \caption{Resultados da seleção de atributos por modalidade}
    \label{tab:selection_summary}
    \begin{tabular}{ccccccc}
        \toprule
        \textbf{Modalidade} & \makecell[c]{\textbf{PCs} \\ \textbf{entrada}} & \makecell[c]{\textbf{Sig.} \\ \textbf{Bonf.}} & \makecell[c]{\textbf{RFE} \\ \textbf{ótimo}} & \makecell[c]{\textbf{L1} \\ \textbf{não-zero}} & \textbf{Consenso} & \textbf{Acc. CV} \\
        \midrule
        EEG & 38 & 2  & 36 & 38 & 36 & 65{,}6\% \\
        ECG & 22 & 0  & 10 & 19 & 10 & 36{,}7\% \\
        EMG & 18 & 0  & 16 & 15 & 16 & 53{,}9\% \\
        \bottomrule
    \end{tabular}
\end{table}

\subsection{EEG}

Para o EEG, o RFECV determinou que o subconjunto ótimo contém 36 dos 38 componentes, atingindo acurácia de validação cruzada de 65{,}6\% para três classes (chance = 33{,}3\%). Apenas 2 componentes sobreviveram à correção de Bonferroni (PC29, $F = 16{,}53$; PC6, $F = 9{,}05$), refletindo a elevada multiplicidade dos 38 testes simultâneos. O componente de maior destaque é o PC29, que apresenta $f^2 = 0{,}089$ (efeito pequeno a médio) e informação mútua de 0{,}117. O fato de o RFECV reter 36 dos 38 componentes sugere que a informação de fase está distribuída difusamente pelo espaço PCA do EEG, sem concentração em poucos componentes dominantes.

\begin{figure}[htbp]
    \centering
    \includegraphics[width=0.95\textwidth]{../output/stage9_selection/figures/anova_fscores_eeg.png}
    \caption{$F$-estatísticas ANOVA por componente principal para EEG.}
    \label{fig:anova_eeg}
\end{figure}

\begin{figure}[htbp]
    \centering
    \includegraphics[width=0.95\textwidth]{../output/stage9_selection/figures/rfe_cv_eeg.png}
    \caption{Curva RFECV para EEG: acurácia de validação cruzada por número de componentes retidos.}
    \label{fig:rfe_eeg}
\end{figure}

\subsection{ECG e EMG}

Para o ECG, nenhum componente atingiu significância estatística após correção de Bonferroni, e a acurácia RFECV de 36{,}7\% é essencialmente ao nível do acaso, corroborando os resultados dos Estágios 7 e 8 que indicam ausência de discriminabilidade de fase no sinal cardíaco com janelas de 5 s. O subconjunto consenso retém 10 dos 22 PCs com base no critério RFE.

O EMG apresenta acurácia moderada de 53{,}9\%, com o componente PC15 destacando-se como o mais informativo: $f^2 = 0{,}260$ (equivalente a $\eta^2 = 0{,}207$ pela relação $\eta^2 = f^2/(1+f^2)$; efeito médio, maior valor entre todas as modalidades), consistente com a hipótese de que a atividade muscular periférica reflete em algum grau as fases de estimulação emocional. Nenhum componente de EMG ou ECG sobreviveu à correção de Bonferroni.

\begin{figure}[htbp]
    \centering
    \includegraphics[width=0.95\textwidth]{../output/stage9_selection/figures/consensus_ranking_eeg.png}
    \caption{Ranking de consenso para EEG: posição dos 15 melhores componentes nos métodos ANOVA, MI, L1-coeficiente e RFE.}
    \label{fig:consensus_eeg}
\end{figure}

\section{Validação Estatística Final}
\label{sec:resultados_validacao}

O décimo estágio verificou a prontidão do conjunto de dados conforme descrito na Seção~\ref{sec:metodo_validacao}. A Tabela~\ref{tab:validation_summary} apresenta os resultados consolidados por modalidade.

\begin{table}[htbp]
    \centering
    \caption{Resultados da validação estatística final por modalidade}
    \label{tab:validation_summary}
    \begin{tabular}{lcccccc}
        \toprule
        \textbf{Modalidade} & \textbf{Obs.} & \textbf{Atrib.} & \textbf{VIF máx.} & \textbf{Raz. imbalan.} & \textbf{$\eta^2$ (top PC)} & \textbf{Pronto} \\
        \midrule
        EEG & 375 & 36 & 1{,}000 & 1{,}00 & 0{,}082 & \checkmark \\
        ECG & 41  & 10 & 1{,}000 & 1{,}08 & 0{,}163 & \checkmark \\
        EMG & 37  & 16 & 1{,}000 & 1{,}27 & 0{,}207 & \checkmark \\
        \bottomrule
    \end{tabular}
\end{table}

\subsection{Multicolinearidade}

O VIF de todos os componentes principais é exatamente $1{,}000$ para as três modalidades, confirmando a ortogonalidade matemática garantida pela PCA. Nenhum componente apresenta redundância linear com os demais, validando a integridade da etapa de redução de dimensionalidade.

\subsection{Separabilidade de Fase}

A Figura~\ref{fig:density_eeg} ilustra as curvas de densidade por fase para os três primeiros componentes do EEG. O componente PC29 apresenta o maior poder discriminante individual ($F = 16{,}53$, $\eta^2 = 0{,}082$), sendo o único componente EEG significativo após correção de Bonferroni. O traço de Pillai multivariado do EEG é de $0{,}516$, indicando separabilidade moderada entre as três fases.

O EMG apresenta o maior $\eta^2$ por componente ($\eta^2 = 0{,}207$ para PC15) e a maior distância de Bhattacharyya entre fases ($d_B = 0{,}376$ entre linha de base e recuperação), sugerindo que a atividade muscular periférica reflete diferenças estruturais entre as condições experimentais. O ECG permanece com separabilidade limitada ($F_{\max} = 3{,}71$, $d_B \leq 0{,}258$), consistente com os resultados dos Estágios 7--9.

\begin{figure}[htbp]
    \centering
    \includegraphics[width=\textwidth]{../output/stage10_validation/figures/density_eeg.png}
    \caption{Curvas de densidade KDE por fase experimental para os três componentes EEG mais discriminativos (PC29, PC13, PC6).}
    \label{fig:density_eeg}
\end{figure}

\subsection{Balanceamento de Classes}

O EEG apresenta balanceamento perfeito entre fases (125 observações cada, razão $= 1{,}00$). O ECG possui razão de $1{,}08$ e o EMG de $1{,}27$, ambos dentro do limiar de $1{,}5$ adotado. Nenhuma modalidade requer reamostragem para equilibrar classes.

\begin{figure}[htbp]
    \centering
    \includegraphics[width=0.85\textwidth]{../output/stage10_validation/figures/class_balance_eeg.png}
    \caption{Balanceamento de classes por fase experimental para o EEG.}
    \label{fig:balance_eeg}
\end{figure}

\subsection{Dataset Final}

O dataset cruzado entre modalidades foi construído pela junção interna nas chaves \texttt{subject\_id} e \texttt{phase}, resultando em 13 sujeitos comuns às três modalidades, 30 observações e 62 atributos ($36_{\text{EEG}} + 10_{\text{ECG}} + 16_{\text{EMG}}$). Os 3 sujeitos excluídos da junção correspondem àqueles sem dados de ECG (s013) ou EMG (s002, s015) desde o Estágio 5. O conjunto EEG isolado (375 observações, 36 atributos) permanece disponível para experimentos de classificação unimodal.

O pipeline de 10 estágios está completo. O dataset final está disponível em \texttt{output/stage10\_validation/data/dataset\_final.csv}, pronto para aplicação de classificadores como SVM, $k$-NN, \textit{Random Forest} e redes neurais.
